% ============================================================
%  SECTION 2: INTERNATIONAL REGULATORY FRAMEWORKS
% ============================================================

\section{International Regulatory Frameworks: What Standards Actually Prohibit}

The assessment of honey quality and authenticity is based on a network
of interrelated normative documents of both global and regional scope.
Understanding their scope and limitations is crucial for evaluating
whether the presence of osmophilic yeasts can constitute an independent
basis for disqualifying a product from commercial trade.

\subsection{Codex Alimentarius, CODEX STAN 12-1981}

The primary document regulating honey quality at the international
level is the Codex Alimentarius Standard for Honey
(CODEX STAN 12-1981, last revision 2023), developed jointly
by FAO and WHO~\cite{codex_stan12}. This standard defines honey
as a natural product made by honey bees from the nectar of plants
or from secretions of living parts of plants or excretions of
plant-sucking insects on the living parts of plants, which the bees
collect, transform by combining with specific substances of their own,
deposit, dehydrate, store, and leave in the honeycomb to ripen and
mature~\cite{codex_stan12}.

In terms of physicochemical parameters, CODEX STAN 12-1981 establishes
the following quantitative limits: water content $\leq 20{,}0\%$
(for heather honey \textit{Calluna} $\leq 23{,}0\%$), content
of~hydroxymethylfurfural (HMF) $\leq 40$~mg/kg (or $\leq 80$~mg/kg
for honeys from countries with a tropical climate), diastase
activity $\geq 8$~Schade units (or $\geq 3$~Schade units for honeys
with naturally low enzymatic activity, such as acacia honey
\textit{Robinia pseudoacacia}), and reducing sugar
content~\cite{codex_stan12}. The standard contains \textbf{no limit
on yeast count expressed in CFU/g}. The only microbiological criterion
is the prohibition on trading honey \emph{``that has begun to ferment
or effervesce''}~\cite{codex_stan12}. This is a distinction of
fundamental legal importance: the disqualifying criterion is
\emph{the state of active fermentation}, not \emph{the potential
cause} (the presence of microorganisms capable of fermentation).

\subsection{EU Directive 2001/110/EC and national law}

At the European Union level, the primary legal act regulating honey
trade is Council Directive 2001/110/EC of 20 December 2001 on honey,
implemented in the Polish legal order by a regulation of the Minister
of Agriculture and Rural Development~\cite{eudir2001110}.
This Directive, modeled on the Codex standard, does not establish
a~CFU/g limit for yeasts or fungi in honey. Microbiological parameters
of honey are subject to general food hygiene provisions, in particular
Regulation (EC) No.~852/2004 of the European Parliament and of the
Council on the hygiene of foodstuffs~\cite{eudir2001110}.

In the control practice of veterinary and sanitary inspection
laboratories, there is an unofficial hygiene threshold of
$100$~CFU/g for the total count of fungi and yeasts, used as a general
indicator of production process hygiene; however, this threshold is
not grounded in specific provisions regarding honey and cannot
constitute an independent basis for declaring the product
non-compliant with the requirements of food law within the meaning
of Article~14 of Regulation~(EC)~No.~178/2002~\cite{eudir2001110}.
Directive 2001/110/EC also permits the possibility of filtering,
heating (to specified temperatures), and blending batches of honey
as legal technological treatments, provided they do not lead to
a~change in its fundamental character~\cite{eudir2001110}.


\subsection{United States Pharmacopeia (USP) and other national standards}

The United States Pharmacopeia (USP) classifies honey as an excipient
for pharmaceutical and food use. USP requirements include
physicochemical parameters similar to those of Codex, but also do not
contain a CFU/g limit for osmophilic yeasts~\cite{tsagkaris2021}.
A similar approach is taken by national standards applicable in the
major honey-exporting countries -- China (GB/T~18796-2012), Argentina
(Código Alimentario Argentino, CAA), and Brazil (IN MAPA No.~89/2000)
-- which define honey quality by chemical parameters and enzymatic
activity, not by yeast count~\cite{tsagkaris2021}.

\subsection{Comparative overview of standards}

Table~\ref{tab:standardy} presents a comparison of key honey quality
parameters in applicable international and regional standards with
respect to microbiological criteria.

\begin{table}[ht]
\begin{adjustwidth}{-\extralength}{0cm}
%} % If the paper is ``preprints'', please uncomment this parenthesis.
%\isPreprints{\begin{tabularx}{\textwidth}{CCCC}}{% This command is only used for ``preprints''.

\centering
\caption{Comparison of honey quality parameters in major international
         standards. Abbreviations: n.a., not applicable (no limit
         in the document); CFU, colony-forming units;
         HMF, hydroxymethylfurfural.}
\label{tab:standardy}
\small
%\begin{tabular}{lcccccc}
\begin{tabularx}{\fulllength}{CCCCCC}
\hline
\textbf{Standard} &
\textbf{Moisture} &
\textbf{HMF} &
\textbf{Diastase} &
\textbf{CFU/g limit} &
\textbf{Criterion} \\
 & (\%) & (mg/kg) & (Schade u.) & \textbf{yeasts} & \textbf{microbiolog.} \\
\hline
Codex STAN 12-1981       & $\leq 20{,}0$ & $\leq 40$  & $\geq 8$   & \textbf{none} & prohibition of fermentation \\
EU Directive 2001/110/EC & $\leq 20{,}0$ & $\leq 40$  & $\geq 8$   & \textbf{none} & prohibition of fermentation \\
USP (honey monograph)    & $\leq 21{,}0$ & n.a.       & n.a.        & \textbf{none} & prohibition of fermentation \\
Reg. (EC) 852/2004       & n.a.          & n.a.       & n.a.        & $\leq 100$    & general hygiene$^{*}$ \\
\hline
\multicolumn{6}{l}{$^{*}$~Limit of 100~CFU/g (fungi + yeasts combined) is of general hygiene character} \\
\multicolumn{6}{l}{\phantom{$^{*}$~}and is not a honey-specific criterion.}
\end{tabularx}
\end{adjustwidth}
\end{table}

\subsection{Regulatory conclusion}

Analysis of the above documents leads to an unequivocal conclusion:
no applicable national or international standard treats \emph{the
presence of osmophilic yeasts per~se} as a basis for disqualifying
honey from commercial trade. The disqualifying criterion is exclusively
the \emph{biological outcome} (active fermentation, effervescence),
not the \emph{potential cause} (the presence of microorganisms capable
of fermentation). This distinction is of fundamental practical
importance: honey containing osmophilic yeasts in any number, but
showing no signs of active fermentation and meeting the moisture
criterion of $\leq 20\%$, complies with the requirements of all
applicable standards~\cite{codex_stan12,eudir2001110}.
A CFU/g measurement alone, without a simultaneous assessment of water
activity ($a_w$), temperature, and storage time, is methodologically
insufficient for declaring the product non-compliant with the
requirements of food law~\cite{tsagkaris2021,schoder2026}.
